\documentclass{article}
\usepackage{graphicx} % Required for inserting images

\title{Babooshka}
\author{Khaled Al-Mahdy}
\date{December 2025}

\begin{document}

`\maketitle

\subsection{Introduction}

\end{document}
\documentclass[12pt]{article}

% ---------- PACKAGES ----------
\usepackage[a4paper,margin=1in]{geometry}
\usepackage{setspace}
\usepackage{times}
\usepackage{titlesec}
\usepackage{hyperref}
\usepackage{graphicx}
\usepackage{microtype}
\usepackage{csquotes}

% ---------- SPACING ----------
\onehalfspacing

% ---------- SECTION STYLE ----------
\titleformat{\section}
  {\normalfont\Large\bfseries}{\thesection}{1em}{}

\titleformat{\subsection}
  {\normalfont\large\itshape}{\thesubsection}{1em}{}

% ---------- TITLE ----------
\title{\textbf{Resolving the Hard Problem of Consciousness:\\
A Relativistic Framework for Measurement and Emergence}}
\author{Khaled Al-Mahdy}
\date{}

\begin{document}

\maketitle
\vspace{1em}

\begin{abstract}
This paper presents a unified theoretical framework addressing the hard problem of consciousness through relativistic, entropy-based, and dynamical models. By reframing consciousness as a measurable, frame-dependent process, the proposed theories dissolve the apparent explanatory gap between physical systems and phenomenal experience while enabling empirical and computational applications.
\end{abstract}

\section{Introduction}
The hard problem of consciousness, as articulated by Chalmers (1995), concerns the apparent inability of physical explanations to account for subjective experience. While functional processes admit mechanistic descriptions, phenomenal qualities appear irreducible. This paper argues that the gap arises from non-relativistic assumptions rather than ontological deficiency.

\section{Relative Consciousness Theory}
Relative Consciousness Theory (CRT) defines consciousness as a ratio:
\[
C_r = \frac{A \times I \times S}{E}
\]
where awareness, integration, and self-awareness scale relative to entropy. Consciousness is thus an internally measured property of system coherence rather than an absolute substance.

\section{Entropy Gradients and Dual Streams}
The Entropy-Gradient Consciousness Model locates optimal cognition at the boundary between order and disorder, while the Degenerative–Generative Stream framework explains how growth and decay may coexist within the same system.

\section{Collective and Human–AI Consciousness}
Synchronized systems exhibit multiplicative amplification of consciousness, extending phenomenology beyond the individual. Artificial intelligence functions as an entropy regulator, enhancing human consciousness without requiring sentience.

\section{Conscious Frame Formalism}
By defining consciousness as trajectories through state space rather than static states, the Conscious Frame Formalism renders subjective experience computable, testable, and simulatable.

\section{Conclusion}
Together, these frameworks offer a coherent, measurable, and interdisciplinary solution to the hard problem of consciousness, repositioning it as a perspectival phenomenon rather than a metaphysical impasse.

\end{document}



\documentclass[12pt]{article}

% ---------- PACKAGES ----------
\usepackage[a4paper,margin=1in]{geometry}
\usepackage{setspace}
\usepackage{times}
\usepackage{titlesec}
\usepackage{hyperref}
\usepackage{graphicx}
\usepackage{microtype}
\usepackage{csquotes}

% ---------- SPACING ----------
\onehalfspacing

% ---------- SECTION STYLE ----------
\titleformat{\section}
  {\normalfont\Large\bfseries}{\thesection}{1em}{}

\titleformat{\subsection}
  {\normalfont\large\itshape}{\thesubsection}{1em}{}

% ---------- TITLE ----------
\title{\textbf{Resolving the Hard Problem of Consciousness:\\
A Relativistic Framework for Measurement and Emergence}}
\author{Khaled Al-Mahdy}
\date{}

\begin{document}

\maketitle
\vspace{1em}

\begin{abstract}
This paper presents a unified theoretical framework addressing the hard problem of consciousness through relativistic, entropy-based, and dynamical models. By reframing consciousness as a measurable, frame-dependent process, the proposed theories dissolve the apparent explanatory gap between physical systems and phenomenal experience while enabling empirical and computational applications.
\end{abstract}

\section{Introduction}
The hard problem of consciousness, as articulated by Chalmers (1995), concerns the apparent inability of physical explanations to account for subjective experience. While functional processes admit mechanistic descriptions, phenomenal qualities appear irreducible. This paper argues that the gap arises from non-relativistic assumptions rather than ontological deficiency.

\section{Relative Consciousness Theory}
Relative Consciousness Theory (CRT) defines consciousness as a ratio:
\[
C_r = \frac{A \times I \times S}{E}
\]
where awareness, integration, and self-awareness scale relative to entropy. Consciousness is thus an internally measured property of system coherence rather than an absolute substance.

\section{Entropy Gradients and Dual Streams}
The Entropy-Gradient Consciousness Model locates optimal cognition at the boundary between order and disorder, while the Degenerative–Generative Stream framework explains how growth and decay may coexist within the same system.

\section{Collective and Human–AI Consciousness}
Synchronized systems exhibit multiplicative amplification of consciousness, extending phenomenology beyond the individual. Artificial intelligence functions as an entropy regulator, enhancing human consciousness without requiring sentience.

\section{Conscious Frame Formalism}
By defining consciousness as trajectories through state space rather than static states, the Conscious Frame Formalism renders subjective experience computable, testable, and simulatable.

\section{Conclusion}
Together, these frameworks offer a coherent, measurable, and interdisciplinary solution to the hard problem of consciousness, repositioning it as a perspectival phenomenon rather than a metaphysical impasse.

\end{document}
\documentclass{article}
\begin{document}
Working.
\end{document}

\subsection{\textbf{\textit{\ref{\cite{}}}}}